% =====================================================================
%  MARL Lecture 2: Solution Concepts I
%  Book: Albrecht, Christianos & Schaefer (2024), Sections 4.1-4.5
%  Compile with:  latexmk -pdf lecture02.tex
% =====================================================================
\documentclass[aspectratio=169,11pt]{beamer}
\usepackage{marl-style}

\title{Multi-Agent Reinforcement Learning}
\subtitle{Lecture 2: Solution Concepts I --- Best Response, Minimax, Nash}
\author{[Instructor name]}
\institute{[Course number]}
\date{[Date]}

\begin{document}

\begin{frame}
  \titlepage
  \centering\small Book: Sections 4.1--4.5
\end{frame}

\begin{frame}{Today}
  \begin{enumerate}
    \item Joint policies and expected returns in any game model \hfill \S4.1
    \item Best responses \hfill \S4.2
    \item Minimax solutions of two-agent zero-sum games \hfill \S4.3
    \item Nash equilibrium \hfill \S4.4
    \item $\epsilon$-Nash equilibrium \hfill \S4.5
  \end{enumerate}
  \bigskip
  \centering
  \begin{tikzpicture}[node distance=5mm]
    \node[bb] (gm) {Game model\\ \scriptsize Lecture 1};
    \node[right=of gm] (plus) {$+$};
    \node[bo,right=of plus] (sc) {Solution concept\\ \scriptsize today and Lecture 3};
    \node[right=of sc] (eq) {$=$};
    \node[bg,right=of eq] (p) {MARL problem};
  \end{tikzpicture}
  \reading{Sections 4.1--4.5. Next lecture: 4.6--4.11.}
\end{frame}

\begin{frame}{What is a solution?}
  \begin{itemize}
    \item In an MDP, a policy is optimal if it maximizes expected return in every state. All optimal policies give the same return.
    \item In a game, agent $i$'s return depends on everyone's policies. ``Optimal'' must be defined relative to the others.
    \item A \textbf{solution} is a \textbf{joint policy} $\pi = (\pi_1, \dots, \pi_n)$ whose expected returns $U_i(\pi)$ satisfy some conditions.
    \item Common-reward games have an obvious candidate: maximize the shared return. With different rewards there is no single obvious answer.
  \end{itemize}
  \bigskip
  Questions game theory asks about every solution concept:
  \begin{enumerate}
    \item \textbf{Existence:} does a solution always exist?
    \item \textbf{Uniqueness:} is there one, several, or infinitely many?
    \item \textbf{Computation and learning:} can we find one efficiently, and do learning rules converge to one?
  \end{enumerate}
\end{frame}

% =====================================================================
\section{Joint policy and expected return}

\begin{frame}{Full histories}
  We define $U_i(\pi)$ once, for POSGs, so it covers stochastic and normal-form games too.
  \bigskip

  \begin{itemize}
    \item \textbf{Full history} up to time $t$: states, joint observations and joint actions,
      \[
        \hh^t = \big( (s^\tau, o^\tau, a^\tau)_{\tau=0}^{t-1},\ s^t,\ o^t \big).
      \]
    \item \textbf{Observation history} extracted from it: $\sigma(\hh^t) = (o^0, \dots, o^t)$, abbreviated $h^t$.
    \item Joint observation probability: $\mathcal{O}(o^t \mid a^{t-1}, s^t) = \prod_{i \in I} \mathcal{O}_i(o_i^t \mid a^{t-1}, s^t)$.
    \item Discounted returns with absorbing states: after termination, the game stays in one state with reward 0 forever. This covers finite and infinite episodes, including one-shot normal-form games.
  \end{itemize}
\end{frame}

\begin{frame}{History-based expected return}
  Enumerate every full history and weight its return by its probability:
  \[
    U_i(\pi) = \lim_{t \to \infty} \E_{\hh^t \sim (\mu, T, \mathcal{O}, \pi)} \big[ u_i(\hh^t) \big]
    = \sum_{\hh^t \in \hat H} \Pr(\hh^t \mid \pi)\, u_i(\hh^t),
  \]
  with
  \begin{align*}
    \Pr(\hh^t \mid \pi) &= \mu(s^0)\, \mathcal{O}(o^0 \mid \emptyset, s^0) \prod_{\tau=0}^{t-1} \pi(a^\tau \mid h^\tau)\, T(s^{\tau+1} \mid s^\tau, a^\tau)\, \mathcal{O}(o^{\tau+1} \mid a^\tau, s^{\tau+1}), \\
    u_i(\hh^t) &= \sum_{\tau=0}^{t-1} \gamma^\tau R_i(s^\tau, a^\tau, s^{\tau+1}).
  \end{align*}
  If agents act \textbf{independently} (book Eq.\ 4.5):
  \[
    \pi(a^\tau \mid h^\tau) = \prod_{j \in I} \pi_j(a_j^\tau \mid h_j^\tau).
  \]
  Correlated equilibrium (Lecture 3) and central learning (Lecture 4) drop this product form.
\end{frame}

\begin{frame}{Recursive expected return}
  Bellman-style recursion over full histories ($s(\hh)$ = last state, $\langle \cdot \rangle$ = concatenation):
  \begin{align*}
    V_i^\pi(\hh) &= \sum_{a \in A} \pi(a \mid \sigma(\hh))\, Q_i^\pi(\hh, a), \\
    Q_i^\pi(\hh, a) &= \sum_{s' \in S} T(s' \mid s(\hh), a) \Big[ R_i(s(\hh), a, s') + \gamma \sum_{o' \in O} \mathcal{O}(o' \mid a, s')\, V_i^\pi(\langle \hh, a, s', o' \rangle) \Big],
  \end{align*}
  and the expected return from the start of the game is
  \[
    U_i(\pi) = \E_{s^0 \sim \mu,\ o^0 \sim \mathcal{O}(\cdot \mid \emptyset, s^0)} \big[ V_i^\pi(\langle s^0, o^0 \rangle) \big].
  \]
  \begin{itemize}
    \item $V_i^\pi(\hh)$: agent $i$'s value when all follow $\pi$ after $\hh$.
    \item $Q_i^\pi(\hh, a)$: value of joint action $a$ after $\hh$, then following $\pi$.
    \item This form is what value iteration for games uses (Lecture 5).
  \end{itemize}
\end{frame}

\begin{frame}{The two definitions are the same tree}
  \centering
  \begin{tikzpicture}[x=1cm,y=1cm,
      dec/.style={circle,draw=cblue,fill=cblue!15,minimum size=6mm,inner sep=0,font=\scriptsize},
      chn/.style={diamond,draw=corange,fill=corange!15,minimum size=5mm,inner sep=0},
      leaf/.style={font=\scriptsize}]
    \node[dec] (r) at (0,0) {$\hh^0$};
    \node[chn] (c1) at (-3.2,-1.4) {};
    \node[chn] (c2) at (3.2,-1.4) {};
    \draw[->,thick,cblue] (r) -- node[above left,font=\scriptsize] {$a$, prob.\ $\pi(a \mid h^0)$} (c1);
    \draw[->,thick,cblue] (r) -- node[above right,font=\scriptsize] {$a'$} (c2);
    \foreach \p/\x [count=\k] in {c1/-4.6,c1/-1.8,c2/1.8,c2/4.6}{
      \node[dec] (n\k) at (\x,-2.9) {$\hh^1$};
      \draw[->,thick,corange] (\p) -- (n\k);
      \draw[->,cgray] (n\k) -- ++(-0.5,-0.9);
      \draw[->,cgray] (n\k) -- ++(0.5,-0.9);
      \node[leaf,cgray] at (\x,-4.1) {$\cdots$};
    }
    \node[font=\scriptsize,corange,anchor=east] at (-3.9,-2.0) {$s', o'$ w.p.\ $T \cdot \mathcal{O}$};
    \node[note,text width=3.4cm,anchor=west] at (5.4,0) {\textcolor{cblue}{circles}: agents choose (policy)\\ \textcolor{corange}{diamonds}: nature chooses (dynamics)};
    \node[note,text width=3.4cm,anchor=west] at (5.4,-3.3) {each root-to-leaf path is one full history $\hh^t$};
  \end{tikzpicture}
  \medskip

  \small History-based: sum over all paths, weighting each path's return $u_i$ by its probability. Recursive: compute $V_i^\pi$ bottom-up through the same tree. Both give $U_i(\pi)$.
\end{frame}

\begin{frame}{Special case 1: normal-form games}
  One step, no states. A joint policy is $\pi = (\pi_1, \dots, \pi_n)$ with $\pi_i \in \Delta(A_i)$, and
  \[
    U_i(\pi) = \sum_{a \in A} \Big( \prod_{j \in I} \pi_j(a_j) \Big) R_i(a).
  \]
  \textbf{Two agents, matrix form.} Let $x = \pi_1$, $y = \pi_2$ and $[R_1]_{kl} = R_1(k, l)$. Then
  \[
    U_1(x, y) = \sum_{k, l} x_k\, y_l\, [R_1]_{kl} = x^\top R_1\, y .
  \]
  \textbf{Key property: linearity in one's own policy.} For fixed $\pi_{-i}$,
  \[
    U_i(\pi_i, \pi_{-i}) = \sum_{a_i \in A_i} \pi_i(a_i)\, U_i(a_i, \pi_{-i}),
    \qquad
    U_i(a_i, \pi_{-i}) = \sum_{a_{-i}} \Big( \prod_{j \neq i} \pi_j(a_j) \Big) R_i(a_i, a_{-i}).
  \]
  Most results in this lecture follow from this one line.
\end{frame}

\begin{frame}{Special case 2: stochastic games with Markov policies}
  With full observability, take stationary policies $\pi_i(a_i \mid s)$. The recursion over histories collapses to states:
  \begin{align*}
    Q_i^\pi(s, a) &= \sum_{s'} T(s' \mid s, a) \big[ R_i(s, a, s') + \gamma V_i^\pi(s') \big], \\
    V_i^\pi(s) &= \sum_{a} \Big( \prod_{j} \pi_j(a_j \mid s) \Big) Q_i^\pi(s, a),
    \qquad U_i(\pi) = \sum_s \mu(s)\, V_i^\pi(s).
  \end{align*}
  \begin{itemize}
    \item $Q_i^\pi$ is defined over \textbf{joint} actions: $\prod_j |A_j|$ entries per state.
    \item For a fixed $\pi$, each $V_i^\pi$ solves a linear system, exactly as in MDP policy evaluation.
    \item $V_i^\pi$ depends on $\pi_{-i}$. There is no optimal policy for agent $i$ in isolation.
  \end{itemize}
\end{frame}

% =====================================================================
\section{Best response}

\begin{frame}{Best response}
  \begin{block}{Book Eq.\ 4.9}
    Given the other agents' policies $\pi_{-i} = (\pi_1, \dots, \pi_{i-1}, \pi_{i+1}, \dots, \pi_n)$,
    \[
      \BR_i(\pi_{-i}) = \argmax_{\pi_i} U_i(\langle \pi_i, \pi_{-i} \rangle).
    \]
  \end{block}
  \begin{itemize}
    \item A \emph{set}: the best response need not be unique.
    \item Most solution concepts today are defined through best responses.
    \item Best-response operators also drive algorithms: fictitious play (\S6.3.1) and joint-action learning with agent models (\S6.3.2), Lecture 6; PSRO (\S9.9), Lecture 16.
  \end{itemize}
\end{frame}

\begin{frame}{Best responses in normal-form games}
  \textbf{Claim.} A deterministic best response always exists, and
  \[
    \pi_i \in \BR_i(\pi_{-i}) \iff \supp(\pi_i) \subseteq \argmax_{a_i \in A_i} U_i(a_i, \pi_{-i}).
  \]
  \textbf{Proof.} Let $M = \max_{a_i} U_i(a_i, \pi_{-i})$. By linearity,
  \[
    U_i(\pi_i, \pi_{-i}) = \sum_{a_i} \pi_i(a_i)\, U_i(a_i, \pi_{-i}) \le \sum_{a_i} \pi_i(a_i)\, M = M,
  \]
  with equality iff $\pi_i(a_i) > 0$ only for actions with $U_i(a_i, \pi_{-i}) = M$. Any maximizing action, played deterministically, attains $M$. \hfill $\square$
  \bigskip

  \begin{itemize}
    \item Computing a best response is a maximization over $|A_i|$ numbers.
    \item Every mixture of best-response actions is also a best response: the BR set is a face of the simplex.
  \end{itemize}
\end{frame}

\begin{frame}{Picture: maximizing a linear function on the simplex}
  \begin{columns}[c]
    \begin{column}{0.45\textwidth}
      Rock-Paper-Scissors. Agent 2 plays $\pi_2 = (0.5, 0.3, 0.2)$ over (R, P, S). Agent 1's action values:
      \begin{align*}
        U_1(R, \pi_2) &= 0(0.5) - 1(0.3) + 1(0.2) = -0.1 \\
        U_1(P, \pi_2) &= 1(0.5) + 0(0.3) - 1(0.2) = \phantom{-}0.3 \\
        U_1(S, \pi_2) &= -1(0.5) + 1(0.3) + 0(0.2) = -0.2
      \end{align*}
      $U_1(\cdot, \pi_2)$ is linear on the simplex $\Delta(A_1)$, so its level sets are parallel lines and the maximum sits at a vertex.
      \medskip

      $\BR_1(\pi_2) = \{P\}$, unique here.
    \end{column}
    \begin{column}{0.52\textwidth}
      \centering
      \begin{tikzpicture}[scale=1.25]
        \coordinate (R) at (0,0);
        \coordinate (P) at (4,0);
        \coordinate (S) at (2,3.464);
        \begin{scope}
          \clip (R) -- (P) -- (S) -- cycle;
          % f(x,y) = 0.1 x - 0.0866 y - 0.1 ; level lines f = k
          \foreach \k in {-0.18,-0.12,...,0.3}{
            \pgfmathsetmacro{\mix}{max(0,min(100,(\k+0.2)*200))}
            \pgfmathsetmacro{\xa}{(\k+0.1+0.0866*(-0.5))/0.1}
            \pgfmathsetmacro{\xb}{(\k+0.1+0.0866*4)/0.1}
            \draw[line width=1.1pt,color=corange!\mix!cblue] (\xa,-0.5) -- (\xb,4);
          }
        \end{scope}
        \draw[thick] (R) -- (P) -- (S) -- cycle;
        \node[below left,font=\small] at (R) {R \scriptsize($-0.1$)};
        \node[below right,font=\small] at (P) {P \scriptsize($0.3$)};
        \node[above,font=\small] at (S) {S \scriptsize($-0.2$)};
        \fill[cred] (P) circle (2.5pt);
        \node[font=\scriptsize,cred,anchor=west] at (4.1,0.35) {maximum};
        \node[note,text width=5cm] at (2,-0.75) {level lines of $U_1(\pi_1, \pi_2)$ over $\pi_1 \in \Delta(\{R,P,S\})$; blue low, orange high};
      \end{tikzpicture}
    \end{column}
  \end{columns}
\end{frame}

\begin{frame}{Best responses need not be unique}
  Let agent 1 play uniformly in Rock-Paper-Scissors, $\pi_1 = (\tfrac13, \tfrac13, \tfrac13)$. For any $\pi_2$ (book footnote 6):
  \begin{align*}
    U_2(\pi_1, \pi_2)
    &= \pi_2(R) \tfrac13 \big[ R_2(R,R) + R_2(P,R) + R_2(S,R) \big] \\
    &\quad + \pi_2(P) \tfrac13 \big[ R_2(R,P) + R_2(P,P) + R_2(S,P) \big] \\
    &\quad + \pi_2(S) \tfrac13 \big[ R_2(R,S) + R_2(P,S) + R_2(S,S) \big] \\
    &= \pi_2(R) \tfrac13 (0 - 1 + 1) + \pi_2(P) \tfrac13 (1 + 0 - 1) + \pi_2(S) \tfrac13 (-1 + 1 + 0) = 0.
  \end{align*}
  \begin{itemize}
    \item \emph{Every} $\pi_2$ is a best response to uniform play: $\BR_2(\pi_1) = \Delta(A_2)$, infinitely many.
    \item In picture terms: the linear function is constant, so the whole simplex is optimal.
  \end{itemize}
\end{frame}

\begin{frame}{Best responses in stochastic games}
  Fix $\pi_{-i}$ (stationary). Marginalizing the others' actions gives an MDP for agent $i$:
  \begin{align*}
    \bar T_i(s' \mid s, a_i) &= \sum_{a_{-i}} \pi_{-i}(a_{-i} \mid s)\, T(s' \mid s, a_i, a_{-i}), \\
    \bar R_i(s, a_i, s') &= \frac{\sum_{a_{-i}} \pi_{-i}(a_{-i} \mid s)\, T(s' \mid s, a_i, a_{-i})\, R_i(s, a_i, a_{-i}, s')}{\bar T_i(s' \mid s, a_i)} .
  \end{align*}
  \begin{itemize}
    \item $\BR_i(\pi_{-i})$ = the optimal policies of the MDP $(S, A_i, \bar T_i, \bar R_i, \gamma)$; a deterministic stationary one exists.
    \item Computing a best response is \textbf{single-agent RL} against fixed opponents.
    \item The catch: $\bar T_i$ and $\bar R_i$ depend on $\pi_{-i}$. If the others change, so does agent $i$'s MDP.
  \end{itemize}
\end{frame}

% =====================================================================
\section{Minimax}

\begin{frame}{Minimax solutions}
  \begin{block}{Definition 5 (Minimax solution)}
    In a zero-sum game with two agents, $\pi = (\pi_i, \pi_j)$ is a minimax solution if
    \[
      U_i(\pi) = \max_{\pi_i'} \min_{\pi_j'} U_i(\pi_i', \pi_j') = \min_{\pi_j'} \max_{\pi_i'} U_i(\pi_i', \pi_j') = -U_j(\pi).
    \]
  \end{block}
  \begin{itemize}
    \item $\max_{\pi_i'} \min_{\pi_j'} U_i$: the most agent $i$ can \textbf{guarantee} against any opponent. $\pi_i$ is $i$'s \emph{maxmin} policy.
    \item $\min_{\pi_j'} \max_{\pi_i'} U_i$: the least agent $j$ can \textbf{force} on $i$. $\pi_j$ is $j$'s \emph{minmax} policy against $i$.
    \item At a minimax solution the two coincide: it does not matter who announces their policy first.
    \item Equivalently: $\pi_i \in \BR_i(\pi_j)$ and $\pi_j \in \BR_j(\pi_i)$.
  \end{itemize}
\end{frame}

\begin{frame}{Weak duality: maxmin $\le$ minmax always}
  \textbf{Claim.} For any function $U$ of two arguments,
  \[
    \max_{x} \min_{y} U(x, y) \le \min_{y} \max_{x} U(x, y).
  \]
  \textbf{Proof.} For any fixed $\bar x$ and $\bar y$:
  \[
    \min_{y} U(\bar x, y) \le U(\bar x, \bar y) \le \max_{x} U(x, \bar y).
  \]
  The left side does not depend on $\bar y$ and the right side does not depend on $\bar x$. Maximize the left over $\bar x$ and minimize the right over $\bar y$. \hfill $\square$
  \bigskip

  \textbf{Theorem (von Neumann, 1928).} For finite two-agent zero-sum normal-form games with mixed policies, equality holds. The common value is the \textbf{value of the game}.
  \begin{itemize}
    \item Several minimax solutions may exist, but all give the same value.
    \item It extends to zero-sum stochastic games with finite episodes or discounting (Shapley, 1953).
    \item It can fail with continuous actions (Sion \& Wolfe, 1957).
  \end{itemize}
\end{frame}

\begin{frame}{Minimax via linear programming}
  Agent $i$ chooses $x_{a_i} = \pi_i(a_i)$ to minimize the best expected return agent $j$ can get (book Eqs.\ 4.12--4.15):
  \begin{align*}
    \text{minimize} \quad & U_j^* \\
    \text{subject to} \quad & \sum_{a_i \in A_i} R_j(a_i, a_j)\, x_{a_i} \le U_j^* && \forall a_j \in A_j \\
    & x_{a_i} \ge 0 \quad \forall a_i \in A_i, \qquad \sum_{a_i \in A_i} x_{a_i} = 1 .
  \end{align*}
  \textbf{Why only pure $a_j$ in the constraints?} By linearity, $U_j(x, \pi_j) = \sum_{a_j} \pi_j(a_j) U_j(x, a_j)$, a weighted average of the left-hand sides. If every pure $a_j$ gives at most $U_j^*$, so does every mixture.
  \begin{itemize}
    \item Agent $j$'s policy comes from the same LP with $i$ and $j$ swapped. At the solution, $U_i^* = -U_j^*$.
    \item The two LPs are duals of each other; LP strong duality gives von Neumann's theorem.
    \item Solvable in polynomial time (interior point), and fast in practice (simplex).
  \end{itemize}
\end{frame}

\begin{frame}{Worked example: a $2 \times 2$ zero-sum game}
  \begin{columns}[c]
    \begin{column}{0.5\textwidth}
      \small
      Agent 1's rewards ($R_2 = -R_1$):
      \begin{center}
        \begin{tabular}{c|cc}
          & L & R \\ \hline
          T & $3$ & $-1$ \\
          B & $-2$ & $1$
        \end{tabular}
      \end{center}
      Let $x = \pi_1(T)$. Agent 1's return against each pure column:
      \begin{align*}
        U_1(x, L) &= 3x - 2(1-x) = 5x - 2, \\
        U_1(x, R) &= -x + (1-x) = 1 - 2x .
      \end{align*}
      Agent 1 maximizes the lower envelope $\min\{5x - 2,\ 1 - 2x\}$:
      \[
        5x - 2 = 1 - 2x \ \Rightarrow\ x^* = \tfrac37, \quad v^* = \tfrac17 .
      \]
      Agent 2, with $y = \pi_2(L)$: $3y - (1-y) = -2y + (1-y) \Rightarrow y^* = \tfrac27$, which gives the same value $\tfrac17$.
    \end{column}
    \begin{column}{0.48\textwidth}
      \centering
      \begin{tikzpicture}[x=4.6cm,y=0.85cm]
        \draw[->] (0,0) -- (1.08,0) node[right,font=\small] {$x$};
        \draw[->] (0,-2.3) -- (0,3.4) node[above,font=\small] {$U_1$};
        \foreach \y in {-2,-1,1,2,3} \draw (-0.01,\y) -- (0.01,\y) node[left=2pt,font=\tiny] {$\y$};
        \draw (1,-0.08) -- (1,0.08) node[below=4pt,font=\tiny] {1};
        \draw[thick,cblue] (0,-2) -- (1,3) node[below right,font=\scriptsize] {vs.\ L};
        \draw[thick,corange] (0,1) -- (1,-1) node[right,font=\scriptsize] {vs.\ R};
        \draw[line width=2.6pt,cred,opacity=0.55] (0,-2) -- (3/7,1/7) -- (1,-1);
        \draw[dashed,cgray] (3/7,0.143) -- (3/7,-2.3) node[below,font=\scriptsize,text=black] {$x^* = \tfrac37$};
        \fill[cred] (3/7,1/7) circle (2.5pt);
        \node[font=\scriptsize,cred,anchor=south east] at (0.41,0.3) {$v^* = \tfrac17$};
        \node[note] at (0.5,-3.2) {red: lower envelope $\min_y U_1(x, y)$};
      \end{tikzpicture}
    \end{column}
  \end{columns}
\end{frame}

\begin{frame}{Rock-Paper-Scissors has a unique minimax solution}
  \textbf{Claim.} Uniform play by both agents is the only minimax solution; the value is 0.
  \bigskip

  \textbf{Proof.} Each agent can guarantee 0 by playing uniformly (previous slide), so the value is $v^* = 0$. Take any $\pi_2$. Agent 1's action values are
  \[
    U_1(R, \pi_2) = \pi_2(S) - \pi_2(P), \quad
    U_1(P, \pi_2) = \pi_2(R) - \pi_2(S), \quad
    U_1(S, \pi_2) = \pi_2(P) - \pi_2(R).
  \]
  They sum to 0. If they are not all 0, the largest is positive, so agent 1 earns more than $v^* = 0$ against $\pi_2$ and $\pi_2$ is not a minmax policy. If they are all 0, then $\pi_2(R) = \pi_2(P) = \pi_2(S)$. Same argument for agent 1. \hfill $\square$
  \bigskip

  \textbf{Interchangeability.} In zero-sum games, if $(\pi_i, \pi_j)$ and $(\pi_i', \pi_j')$ are minimax solutions, so are $(\pi_i, \pi_j')$ and $(\pi_i', \pi_j)$. Each agent can compute its part alone, with no coordination. Self-play for zero-sum games (AlphaZero, Lecture 15) relies on this.
\end{frame}

% =====================================================================
\section{Nash equilibrium}

\begin{frame}{Nash equilibrium}
  \begin{block}{Definition 6 (Nash equilibrium; Nash, 1950)}
    In a general-sum game with $n$ agents, $\pi = (\pi_1, \dots, \pi_n)$ is a Nash equilibrium if
    \[
      \forall i,\ \forall \pi_i' : \quad U_i(\pi_i', \pi_{-i}) \le U_i(\pi).
    \]
  \end{block}
  \begin{itemize}
    \item No agent gains by \textbf{unilaterally} changing its policy.
    \item Equivalently: $\pi_i \in \BR_i(\pi_{-i})$ for all $i$, a \emph{mutual} best response.
    \item In two-agent zero-sum games, Nash equilibria and minimax solutions coincide.
    \item In normal-form games it suffices to check deterministic deviations (best-response claim): a finite set of inequalities
      \[
        U_i(a_i, \pi_{-i}) \le U_i(\pi) \qquad \forall i,\ \forall a_i \in A_i .
      \]
  \end{itemize}
\end{frame}

\begin{frame}{Equilibria in the three book games}
  \begin{columns}[T]
    \begin{column}{0.32\textwidth}
      \centering \textbf{Prisoner's Dilemma}\\[3pt]
      \footnotesize
      \begin{tabular}{c|cc}
          & C & D \\ \hline
        C & $-1,-1$ & $-5,0$ \\
        D & $0,-5$ & \cellcolor{cred!20}$-3,-3$
      \end{tabular}\\[5pt]
      \small one NE, deterministic: $(D,D)$
    \end{column}
    \begin{column}{0.32\textwidth}
      \centering \textbf{Coordination}\\[3pt]
      \footnotesize
      \begin{tabular}{c|cc}
          & A & B \\ \hline
        A & \cellcolor{cgreen!20}$10$ & $0$ \\
        B & $0$ & \cellcolor{cgreen!20}$10$
      \end{tabular}\\[5pt]
      \small three NE: $(A,A)$, $(B,B)$, and $(\tfrac12, \tfrac12)$
    \end{column}
    \begin{column}{0.32\textwidth}
      \centering \textbf{Rock-Paper-Scissors}\\[3pt]
      \footnotesize
      \begin{tabular}{c|ccc}
          & R & P & S \\ \hline
        R & $0,0$ & $-1,1$ & $1,-1$ \\
        P & $1,-1$ & $0,0$ & $-1,1$ \\
        S & $-1,1$ & $1,-1$ & $0,0$
      \end{tabular}\\[5pt]
      \small one NE, probabilistic: uniform
    \end{column}
  \end{columns}
  \bigskip
  \begin{itemize}
    \item Shaded cells are deterministic equilibria.
    \item Two lessons: equilibria can be \textbf{deterministic or probabilistic}, and a game can have \textbf{several equilibria with different returns}.
  \end{itemize}
\end{frame}

\begin{frame}{Dominance: why $(D,D)$ in the Prisoner's Dilemma}
  Action $a_i$ \textbf{strictly dominates} $a_i'$ if $R_i(a_i, a_{-i}) > R_i(a_i', a_{-i})$ for every $a_{-i}$.
  \bigskip

  For agent 1:
  \[
    R_1(D, C) = 0 > -1 = R_1(C, C), \qquad R_1(D, D) = -3 > -5 = R_1(C, D).
  \]
  So D strictly dominates C; by symmetry the same for agent 2.
  \begin{itemize}
    \item A strictly dominated action is never in the support of a best response, hence never played in any Nash equilibrium.
    \item The only equilibrium is $(D,D)$ with $(-3,-3)$, although $(C,C)$ gives $(-1,-1)$ to both.
    \item Equilibrium is about \emph{stability}, not \emph{efficiency}. Pareto optimality and welfare (Lecture 3) address efficiency.
  \end{itemize}
\end{frame}

\begin{frame}{The indifference principle}
  \textbf{Corollary of the best-response claim.} At a Nash equilibrium $\pi$, for every agent $i$:
  \[
    U_i(a_i, \pi_{-i}) = U_i(\pi) \quad \forall a_i \in \supp(\pi_i), \qquad
    U_i(a_i, \pi_{-i}) \le U_i(\pi) \quad \forall a_i \in A_i .
  \]
  Every action an agent mixes over earns the same, and nothing earns more.
  \bigskip

  \textbf{Recipe ($2 \times 2$):} pick the \emph{other} agent's mixture so that \emph{you} are indifferent.
  \bigskip

  \textbf{Coordination game.} Let $y = \pi_2(A)$. Agent 1:
  \[
    U_1(A, y) = 10y, \qquad U_1(B, y) = 10(1 - y).
  \]
  Indifference: $10y = 10(1-y) \Rightarrow y = \tfrac12$. By symmetry $\pi_1(A) = \tfrac12$. Expected return:
  \[
    U_i = \tfrac14 \cdot 10 + \tfrac14 \cdot 10 = 5 < 10 .
  \]
  The probabilistic equilibrium is worse for both than either deterministic one.
\end{frame}

\begin{frame}{Stag Hunt: drawing the best-response graphs}
  \begin{columns}[T]
    \begin{column}{0.5\textwidth}
      \centering
      \footnotesize
      \begin{tabular}{c|cc}
          & S & H \\ \hline
        S & $4,4$ & $0,3$ \\
        H & $3,0$ & $3,3$
      \end{tabular}
      \raggedright
      \medskip

      \small
      $p = \pi_1(S)$, $q = \pi_2(S)$. Agent 1:
      \[
        U_1(S, q) = 4q, \qquad U_1(H, q) = 3 .
      \]
      \[
        \BR_1(q) =
        \begin{cases}
          H\ (p = 0) & q < \tfrac34 \\
          \text{any } p & q = \tfrac34 \\
          S\ (p = 1) & q > \tfrac34
        \end{cases}
      \]
      Symmetric for agent 2. Equilibria are where the graphs cross:
      \[
        (p, q) \in \{(0,0),\ (\tfrac34, \tfrac34),\ (1,1)\}.
      \]
      Book \S5.4.2 uses this game for equilibrium selection (Lecture 4).
    \end{column}
    \begin{column}{0.46\textwidth}
      \centering
      \begin{tikzpicture}[scale=4.3]
        \fill[cgray!6] (0,0) rectangle (1,1);
        \draw[->] (0,0) -- (1.1,0) node[right,font=\small] {$q$};
        \draw[->] (0,0) -- (0,1.1) node[above,font=\small] {$p$};
        \foreach \t/\l in {0/0,0.75/\tfrac34,1/1}{
          \draw (\t,0) -- (\t,-0.02) node[below,font=\scriptsize] {$\l$};
          \draw (0,\t) -- (-0.02,\t) node[left,font=\scriptsize] {$\l$};}
        \draw[cblue,line width=1.6pt] (0,0) -- (0.75,0) -- (0.75,1) -- (1,1);
        \draw[corange,line width=1.6pt,dashed] (0,0) -- (0,0.75) -- (1,0.75) -- (1,1);
        \foreach \x/\y in {0/0,0.75/0.75,1/1} \fill[cred] (\x,\y) circle (0.022);
        \node[font=\scriptsize,cred,anchor=south west] at (1,1) {$(S,S)$: 4, 4};
        \node[font=\scriptsize,cred,anchor=north east] at (0.73,0.73) {mixed: 3, 3};
        \node[font=\scriptsize,cred,anchor=north west] at (0.02,0.0) {$(H,H)$: 3, 3};
        \node[font=\scriptsize,cblue,anchor=west] at (0.77,0.4) {$\BR_1$};
        \node[font=\scriptsize,corange,anchor=south] at (0.4,0.76) {$\BR_2$};
      \end{tikzpicture}
    \end{column}
  \end{columns}
\end{frame}

\begin{frame}{Deterministic versus probabilistic equilibria}
  \begin{itemize}
    \item Some games have \textbf{only} probabilistic equilibria (Rock-Paper-Scissors, Matching Pennies).
    \item Algorithms that represent only deterministic policies, such as greedy Q-learning, \textbf{cannot} represent these equilibria, let alone learn them (Lecture 5 onward).
    \item With three or more agents, equilibrium probabilities can be \textbf{irrational} numbers (Nash, 1950): no finite-precision representation. This motivates $\epsilon$-Nash.
  \end{itemize}
  \bigskip
  \begin{center}
    \begin{tabular}{lcc}
      \toprule
      Game & deterministic NE & probabilistic NE \\
      \midrule
      Prisoner's Dilemma & $(D,D)$ & none \\
      Coordination & $(A,A)$, $(B,B)$ & $(\tfrac12, \tfrac12)$ \\
      Stag Hunt & $(S,S)$, $(H,H)$ & $(\tfrac34, \tfrac34)$ \\
      Rock-Paper-Scissors & none & uniform \\
      \bottomrule
    \end{tabular}
  \end{center}
\end{frame}

\begin{frame}{Existence of Nash equilibria}
  \small
  \textbf{Theorem (Nash, 1950).} Every finite normal-form game has at least one Nash equilibrium. Stochastic games too (Fink, 1964).
  \bigskip

  \textbf{Proof sketch.} Gain from switching to deterministic $a$: \ $g_i^a(\pi) = \max\{0,\ U_i(a, \pi_{-i}) - U_i(\pi)\}$. Define
  \[
    f(\pi)_i(a) = \frac{\pi_i(a) + g_i^a(\pi)}{1 + \sum_{b \in A_i} g_i^b(\pi)} .
  \]
  \begin{enumerate}
    \item $f$ maps the product of simplices to itself and is continuous; by Brouwer it has a fixed point $\pi$.
    \item $U_i(\pi)$ is an average of $U_i(a, \pi_{-i})$ over $a \in \supp(\pi_i)$, so some $a$ in the support has $U_i(a, \pi_{-i}) \le U_i(\pi)$, hence $g_i^a(\pi) = 0$.
    \item At that $a$ the fixed-point equation reads $\pi_i(a) = \pi_i(a) / (1 + \sum_b g_i^b)$ with $\pi_i(a) > 0$, so $\sum_b g_i^b(\pi) = 0$.
    \item No deterministic deviation gains for any $i$: $\pi$ is a Nash equilibrium. \hfill $\square$
  \end{enumerate}
  \medskip
  Existence is non-constructive. Computing an equilibrium is hard in general (PPAD, \S4.11, Lecture 3).
\end{frame}

\begin{frame}{Folk theorems: almost anything is an equilibrium}
  \begin{columns}[c]
    \begin{column}{0.54\textwidth}
      \small
      In repeated or stochastic games with $\gamma$ close to 1, any returns $\hat U = (\hat U_1, \dots, \hat U_n)$ that are
      \begin{itemize}
        \item \textbf{feasible}: $U_i(\pi) = \hat U_i$ for some joint policy $\pi$, and
        \item \textbf{enforceable}: $\hat U_i \ge v_i$, agent $i$'s minmax value
          \[
            v_i = \min_{\pi_{-i}^{mm}} \max_{\pi_i^{mm}} U_i(\pi_i^{mm}, \pi_{-i}^{mm}),
          \]
      \end{itemize}
      can be achieved by an equilibrium.
      \medskip

      \textbf{Construction:} play $\pi$; if agent $i$ ever deviates, the others play $\pi_{-i}^{mm}$ forever. Deviating then earns at most $v_i \le \hat U_i$.
      \medskip

      Repeated PD: $v_i = -3$ (the other defects). Mutual cooperation, $(-1,-1)$, is an equilibrium outcome when agents are patient enough.
    \end{column}
    \begin{column}{0.44\textwidth}
      \centering
      \begin{tikzpicture}[scale=0.82]
        \fill[cblue!10] (-5,0) -- (-1,-1) -- (0,-5) -- (-3,-3) -- cycle;
        \fill[cgreen!35] (-3,-3) -- (-3,-0.5) -- (-1,-1) -- (-0.5,-3) -- cycle;
        \draw[->] (-5.6,0) -- (0.7,0) node[right,font=\scriptsize] {$U_1$};
        \draw[->] (0,-5.6) -- (0,0.7) node[above,font=\scriptsize] {$U_2$};
        \foreach \t in {-5,-3,-1} {
          \draw (\t,0.08) -- (\t,-0.08) node[above=3pt,font=\tiny] {$\t$};
          \draw (0.08,\t) -- (-0.08,\t) node[right=3pt,font=\tiny] {$\t$};}
        \draw[dashed,cred] (-3,-5.3) -- (-3,0.3);
        \draw[dashed,cred] (-5.3,-3) -- (0.3,-3);
        \fill[cred] (-3,-3) circle (2.5pt) node[below left,font=\scriptsize,cred] {$(v_1, v_2)$};
        \fill[cgreen!60!black] (-1,-1) circle (2.5pt) node[below left,font=\scriptsize] {$(C,C)$};
        \fill (-5,0) circle (1.8pt) node[below=2pt,font=\scriptsize] {$(C,D)$};
        \fill (0,-5) circle (1.8pt) node[left,font=\scriptsize] {$(D,C)$};
        \node[note,text width=4.6cm] at (-2.5,-6.2) {\textcolor{cblue}{blue}: feasible returns; \textcolor{cgreen}{green}: feasible and enforceable};
      \end{tikzpicture}
    \end{column}
  \end{columns}
\end{frame}

\begin{frame}{Checking whether $\pi$ is a Nash equilibrium}
  The definition reduces one multi-agent check to $n$ single-agent problems:
  \bigskip

  \centering
  \begin{tikzpicture}[font=\small]
    \node[b,text width=2.0cm] (pi) at (0,0) {joint policy $\pi$};
    \node[bb,text width=3.2cm] (br) at (3.3,0) {for each agent $i$:\\ fix $\pi_{-i}$, compute a best response $\pi_i'$};
    \node[bo,text width=3.0cm] (cmp) at (6.85,0) {$U_i(\pi_i', \pi_{-i}) > U_i(\pi)$?};
    \node[br,text width=2.3cm,font=\footnotesize] (no) at (11.1,0.9) {not a Nash equilibrium};
    \node[bg,text width=2.3cm,font=\footnotesize] (yes) at (11.1,-0.9) {Nash equilibrium};
    \draw[->,thick] (pi) -- (br);
    \draw[->,thick] (br) -- (cmp);
    \draw[->,thick,cred] (cmp.east) -- ++(0.3,0) |- node[pos=0.72,above,font=\tiny,text=cred] {yes, some $i$} (no.west);
    \draw[->,thick,cgreen] (cmp.east) -- ++(0.3,0) |- node[pos=0.72,below,font=\tiny,text=cgreen] {no, every $i$} (yes.west);
  \end{tikzpicture}
  \bigskip

  \raggedright
  \begin{itemize}
    \item Normal-form games: each best response is an $\arg\max$ over $|A_i|$ values, or an LP.
    \item Sequential games: each best response is an MDP, solved by dynamic programming or single-agent RL.
    \item The gap $\sum_i [\max_{\pi_i'} U_i(\pi_i', \pi_{-i}) - U_i(\pi)]$ (``NashConv'', exploitability in zero-sum games) measures distance from equilibrium. Used to evaluate PSRO (Lecture 16).
  \end{itemize}
\end{frame}

% =====================================================================
\section{$\epsilon$-Nash equilibrium}

\begin{frame}{$\epsilon$-Nash equilibrium}
  \textbf{Why relax Nash?}
  \begin{itemize}
    \item Equilibrium probabilities may be irrational; computers use finite precision.
    \item Reaching an exact equilibrium can be too expensive. Small remaining gains may be acceptable.
  \end{itemize}
  \begin{block}{Definition 7 ($\epsilon$-Nash equilibrium)}
    For $\epsilon > 0$, $\pi$ is an $\epsilon$-Nash equilibrium if
    \[
      \forall i,\ \forall \pi_i' : \quad U_i(\pi_i', \pi_{-i}) - \epsilon \le U_i(\pi).
    \]
  \end{block}
  \begin{itemize}
    \item Exact Nash is $\epsilon = 0$.
    \item $U_i$ is continuous in $\pi$, so every Nash equilibrium is surrounded by a region of $\epsilon$-Nash equilibria, for every $\epsilon > 0$.
    \item Checking it: the same best-response procedure, testing $U_i(\pi_i', \pi_{-i}) - \epsilon > U_i(\pi)$.
  \end{itemize}
\end{frame}

\begin{frame}{An $\epsilon$-Nash equilibrium need not be near any Nash equilibrium}
  \small
  \vspace{3mm}
  \begin{columns}[T]
    \begin{column}{0.36\textwidth}
      \centering
      \begin{tabular}{c|cc}
          & C & D \\ \hline
        A & \cellcolor{cgreen!25}$100,100$ & $0,0$ \\
        B & $1,2$ & \cellcolor{corange!25}$1,1$
      \end{tabular}\\[4pt]
      \small book Figure 4.2
    \end{column}
    \begin{column}{0.6\textwidth}
      \textbf{$(A,C)$ is the unique Nash equilibrium.}
      \begin{itemize}
        \item For agent 2, C strictly dominates D: $100 > 0$ and $2 > 1$. So agent 2 plays C in every equilibrium.
        \item Against C, agent 1 prefers A ($100 > 1$). Hence $(A,C)$ is the only one.
      \end{itemize}
      \textbf{$(B,D)$ is a $1$-Nash equilibrium.}
      \begin{itemize}
        \item Agent 1 against D: A gives 0, B gives 1. No gain.
        \item Agent 2 against B: C gives 2, D gives 1. Gain $= 1 = \epsilon$.
      \end{itemize}
    \end{column}
  \end{columns}
  \bigskip
  \begin{itemize}
    \item Returns $(1,1)$ versus $(100,100)$: neither agent is within $\epsilon$ of its equilibrium return.
    \item Raise the $(A,C)$ rewards to anything: $(B,D)$ stays a $1$-Nash equilibrium. The gap is unbounded.
    \item Small $\epsilon$ means ``no agent gains much by deviating alone'', not ``close to an equilibrium''.
  \end{itemize}
\end{frame}

% =====================================================================
\section{Summary}

\begin{frame}{Summary}
  \begin{itemize}
    \item $U_i(\pi)$ is defined for every game model, by enumerating histories or by recursion. In normal-form games it is linear in each $\pi_i$.
    \item \textbf{Best response:} $\BR_i(\pi_{-i}) = \argmax_{\pi_i} U_i$. Always has a deterministic member; may be a large set. In stochastic games, it is an MDP.
    \item \textbf{Minimax} (two-agent zero-sum): maxmin = minmax; unique value; interchangeable solutions; computed by LP.
    \item \textbf{Nash:} mutual best response. Always exists; may be probabilistic; often not unique; not necessarily efficient. Folk theorems: many equilibria in repeated games.
    \item \textbf{$\epsilon$-Nash:} bounded incentive to deviate; can be far from any Nash equilibrium.
  \end{itemize}
  \bigskip
  \textbf{Next lecture} (\S4.6--4.11): correlated equilibria, Pareto optimality, welfare and fairness, no-regret, and the complexity of computing equilibria.
\end{frame}

\begin{frame}{Exercises}
  \begin{enumerate}
    \item Show that the history-based and recursive definitions of $U_i(\pi)$ agree for a two-step stochastic game with two states and two actions per agent.
    \item Show that uniform play is the unique Nash equilibrium of Matching Pennies, using the indifference principle.
    \item Find all Nash equilibria of Chicken: $(D,D) = (0,0)$, $(D,Y) = (7,2)$, $(Y,D) = (2,7)$, $(Y,Y) = (6,6)$.
    \item Write the minimax LP for Rock-Paper-Scissors and verify the solution with an LP solver.
    \item Prove interchangeability: if $(x, y)$ and $(x', y')$ are minimax solutions of a two-agent zero-sum game, so is $(x, y')$.
    \item In the repeated Prisoner's Dilemma with continuation probability $\gamma$, for which $\gamma$ is ``cooperate until the other defects, then defect forever'' (grim trigger) a Nash equilibrium?
    \item Construct a game where an $\epsilon$-Nash equilibrium with $\epsilon = 0.01$ gives every agent a return at least 1000 below its return in every Nash equilibrium.
  \end{enumerate}
\end{frame}

\begin{frame}{References}
  \footnotesize
  \begin{itemize}
    \item Albrecht, Christianos, Sch\"afer. \emph{Multi-Agent Reinforcement Learning: Foundations and Modern Approaches}. MIT Press, 2024. Sections 4.1--4.5.
    \item Fink. Equilibrium in a stochastic $n$-person game. \emph{J.\ Sci.\ Hiroshima Univ.}, 1964.
    \item Fudenberg, Levine. \emph{The Theory of Learning in Games}. MIT Press, 1998.
    \item Nash. Equilibrium points in $n$-person games. \emph{PNAS}, 1950.
    \item Owen. \emph{Game Theory}. 4th ed., Emerald, 2013.
    \item Shapley. Stochastic games. \emph{PNAS}, 1953.
    \item Shoham, Leyton-Brown. \emph{Multiagent Systems}. Cambridge University Press, 2008.
    \item Sion, Wolfe. On a game without a value. \emph{Contributions to the Theory of Games III}, 1957.
    \item von Neumann. Zur Theorie der Gesellschaftsspiele. \emph{Math.\ Annalen}, 1928.
    \item von Neumann, Morgenstern. \emph{Theory of Games and Economic Behavior}. Princeton, 1944.
  \end{itemize}
\end{frame}

\end{document}
